\chapter{Herramientas de gestión}

\begin{quotation} [Escritor] {Tom Clancy (1947--2013)}
La tecnología es sólo una herramienta. La gente usa las herramientas para mejorar sus vidas.
\end{quotation}

\begin{abstract}
Para el correcto desarrollo de proyectos se han de utilizar herramientas tanto generales que controlen el tiempo que se le dedica al proyecto como específicas en cuanto a lo que respecta a proyectos informáticos, como repositorios de código, veamos entonces qué herramientas hemos usado.
\end{abstract}

\section{Herramientas de gestión}
\subsection{Herramienta de gestión del tiempo}
Para la gestión del tiempo se ha decidido el uso de una herramienta en linea que almacene qué se le dedica diariamente en qué intervalos de tiempo a qué actividad. Esto no solo nos permitirá tener en control qué se hace en qué momento sino que facilitará el seguimiento de la planificación y si se están obteniendo los resultados que se esperaban, pudiendo así si fuera necesario volver a planificar el proyecto. 

Para llevar a cabo todas estas mediciones se ha decidido utilizar la herramienta de Clockify, debido a que esta es de uso gratuito y el equipo de desarrollo contaba con experiencia en el uso de esta.

\subsection{Herramienta de gestión del código}

Para la gestión del código, se buscaba que esta fuera un repositorio en el que los desarrolladores pudieran subir los avances de código de manera regular para que este no sea vulnerable a pérdidas en el disco duro local del desarrollador. 

En este respecto se ha decidido usar git, mediante la  herramienta gráfica en línea de GitHub, creando un equipo para los dos desarrolladores y dentro del equipo tres repositorios de código uno para cada respectivo Front-End en el que solo tendrán permisos el desarrollador encargado de ese desarrollo el subir cambios al código, así yo solo tendría permisos para añadir código al Front-End referente a la aplicación web, y por último un tercer repositorio de código en el que se añadan todas las funcionalidades del Back-End en el que ambos desarrolladores puedan añadir las funcionalidades asignadas a cada uno de ellos en el transcurso de las iteraciones.

\subsection{Herramienta de despliegue del código}

Una vez se vayan produciendo iteraciones, para el acceso de ambos desarrolladores de los respectivos Front-End a la base de datos y Back-End del sistema se deberá desplegar en linea para que la última versión del Back-End sea accesible por parte de ambos desarrolladores.

Por ello se ha decidido utilizar la herramienta de despliegue en linea de Heroku debido al conocimiento de esta por parte de los desarrolladores, para el despliegue de aplicaciones y debido a que el despliegue en Heroku es gratuito.

\subsection{Herramienta de gestión y asignación de tareas}

Para la gestión y asignación de tareas hemos decidido el uso de la herramienta de GitHub proyects, en el que no solo obtendremos un repositorio en el que añadiremos las tareas a desarrollar, sino que se podrán asignar a desarrolladores cada tarea y la nomenclatura de la tarea será de utilidad en las ramas que se usarán en la herramienta de gestión del código. Además de esto gracias a columnas se podrá ir siguiendo el avance de las tareas en tiempo real. 

\subsection{Herramienta de redacción de documentación}

En cuanto a la herramienta de redacción de documentación se buscaba una herramienta que permitiera la edición de esta por parte de ambos desarrolladores, la compartición de archivos y un guardado en la nube que evitase pérdidas en caso de que no se guarde el documento de manera recurrente. 

Con estas especificaciones en mente quedaron en mente dos herramientas, la primera de estas Microsoft Word en su versión web es capaz de cumplir los requisitos que se buscábamos. Además, la herramienta en línea de Overleaf, basada en LaTeX que además de cumplir los requisitos previamente mencionados, nos otorga facilidades en cuanto a la exportación y maquetación del documento en formato pdf que se presentará.

Todo lo expuesto anteriormente además de las recomendaciones del supervisor del proyecto, resultaron en la decisión final del uso de la segunda de las herramientas como la idónea para la redacción de la documentación.

\subsection{Herramientas de producción de imágenes}

Para la producción de imágenes se han utilizado dos herramientas en función del tipo de imagen que se buscaba producir con estas. 

Primeramente se han de producir imágenes de la identidad corporativa como pueden ser iconos, se ha decidido usar la herramienta de Canva en su versión gratuita debido a la facilidad de uso de la herramienta, así como el conocimiento de uno de los desarrolladores del uso de la herramienta.

En cuanto a la herramienta gráfica destinada a la producción de los modelos de datos a integrar en la presente memoria se han valorado varias de estas aunque finalmente se ha decidido debido a la sencillez, gratuidad y conocimiento de la herramienta por parte de los desarrolladores utilizar Lucidchart.